\documentclass[12pt]{article}
\usepackage{graphicx}
\usepackage{hyperref}
\usepackage{amsmath}
\usepackage{amssymb}
\usepackage{siunitx}
\usepackage[numbers, square, sort&compress]{natbib}

\begin{document}

\title{Chapter 3: Scoring Functions}
\author{Elvis Martis}
\date{}
\maketitle


\section{Thermodynamics of protein--ligand binding}

Small molecules, acting as messengers or substrates, bind to their respective target receptors to exert their biological or biochemical functions. This is a central process in biochemistry and a fundamental concept essential to understanding to excel at drug discovery. From a thermodynamic perspective, binding is a reversible association between a macromolecular receptor, usually a protein, and a small molecule. This process is characterised by equilibrium constants and changes in Gibbs free energy, enthalpy, and entropy\cite{CH3_DillBromberg,CH3_AtkinsDePaula}. For computational chemists in general, and in particular, for the development of scoring functions in molecular docking, it is essential to understand how these thermodynamic quantities are defined, how they relate to experimental observables such as dissociation constants, and how they can be interpreted in terms of microscopic interactions and molecular motions\cite{CH3_GilsonZhou2007}.

\subsection{Macroscopic description of binding equilibria}

Consider a protein P and a ligand L that reversibly form a complex PL:

\begin{equation}
    \textrm{P + L} \rightleftharpoons \textrm{PL}
\label{eq:PL_compl}
\end{equation}

Under conditions where the solution is sufficiently dilute and ideal, this equilibrium can be characterised by an association constant $K_a$ or a dissociation constant $K_d$. Using molar concentrations, the relation can be written as follows:

\begin{equation}
        \textrm{K}_\textrm{a} = \frac{[PL]}{[P][L]}, \qquad \textrm{K}_\textrm{d} = \frac{[P][L]}{[PL]} = \frac{1}{K_a}
\label{eq:molar_binding}
\end{equation}

Because equilibrium constants are dimensionless, a standard state concentration $c^\circ$, typically $c^\circ$ = 1 M is introduced. The thermodynamic standard Gibbs free energy of binding, $\Delta G^\circ$, is related to $\textrm{K}_\textrm{a}$ as shown in equation \ref{eq:dg_ka}

\begin{equation}
    \Delta G^\circ = -RT \ln K_a,
\label{eq:dg_ka}
\end{equation}

where R is the gas constant, and T is the absolute temperature.\cite{CH3_DillBromberg,CH3_AtkinsDePaula} 

A more explicit definition that emphasises the standard state is

\begin{equation}
    K_a = \frac{[PL]/c^\circ}{([P]/c^\circ)([L]/c^\circ)}, \qquad
    \Delta G^\circ = -RT \ln K_a
\label{eq:dg_ka_ext}
\end{equation}

ensuring that the logarithm is dimensionless.

In drug discovery, affinity is reported as $\textrm{K}_\textrm{d}$, or as an inhibitory concentration such as $\textrm{K}_\textrm{i}$ or $\textrm{IC}_{\textrm{50}}$; under suitable assumptions, these quantities can be related to $\Delta G^\circ$\cite{CH3_GilsonZhou2007}. For example, at T = 298 K, a tenfold change in $\textrm{K}_\textrm{d}$ corresponds to a change in $\Delta G^\circ$ of $\approx$ 1.36 kcal/mol, illustrating how small energetic differences translate into large changes in affinity.

\subsection{Enthalpy, entropy, and the van't Hoff relation}

The standard free energy of binding can be decomposed into standard enthalpic and entropic contributions:

\begin{equation}
    \Delta G^\circ = \Delta H^\circ - T \Delta S^\circ.
\label{eq:gibbs_eq}
\end{equation}

Here, $\Delta H^\circ$ is the standard enthalpy change and $\Delta S^\circ$ is the standard entropy change upon binding.\cite{CH3_DillBromberg} Conceptually, $\Delta H^\circ$ reflects changes in the average internal energy of the system (including protein, ligand, and solvent) at constant pressure. At the same time, $\Delta S^\circ$ incorporates changes in the number of accessible microstates, including translational, rotational, and conformational degrees of freedom of the protein, ligand, and surrounding solvent.

Experimentally, $\Delta H^\circ$ and $\Delta S^\circ$ can be obtained from calorimetric measurements, such as isothermal titration calorimetry (ITC), or from the temperature dependence of the equilibrium constant.\cite{CH3_holdgate2005} Under the approximation that $\Delta H^\circ$ is independent of temperature over a limited range, the van 't Hoff equation relates the temperature derivative of $\ln \textrm{K}_\textrm{a}$ to $\Delta H^\circ$:

\begin{equation}
    \left( \frac{\partial \ln K_a}{\partial (1/T)} \right)_p = -\frac{\Delta H^\circ}{R}.
\label{eq:vant_hoff1}
\end{equation}

Integrating between two temperatures $\textrm{T}_\textrm{1}$ and $\textrm{T}_\textrm{2}$ gives

\begin{equation}
    \ln \frac{K_a(T_2)}{K_a(T_1)} = -\frac{\Delta H^\circ}{R} \left( \frac{1}{T_2} - \frac{1}{T_1} \right).
\label{eq:vant_hoff2}
\end{equation}

The enthalpy and entropy of binding often exhibit compensation: more favourable (negative) $\Delta H^\circ$ is accompanied by less favourable (more negative) $\Delta S^\circ$, and vice versa.\cite{CH3_reynolds2011} This enthalpy--entropy compensation is a common feature of protein--ligand systems and complicates intuitive reasoning about the driving forces of binding.

\subsection{Thermodynamics of covalent binding}

For covalent inhibitors\cite{CH3_Strelow2017,CH3_sutanto2020}, binding is typically described as a two-step process:

\begin{equation}
    \textrm{P + L} \;\underset{k_{\!off}}{\overset{k_{\!on}}{\rightleftharpoons}}\; P\!:\!L
    \;\xrightarrow{k_{\textrm{inact}}}\;
    P\!-\!L,
\label{eq:covbind_eq1}
\end{equation}

where P is the protein, L the ligand, $P\!:\!L$ a reversible noncovalent complex, and $P\!-\!L$ the covalently modified protein. The first step is governed by classical binding thermodynamics: formation of $P\!:\!L$ is characterised by an equilibrium constant

\begin{equation}
    \textrm{K}_\textrm{I} = \frac{k_{\textrm{off}}}{k_{\textrm{on}}} 
\label{covbind_eq2a}
\end{equation}

and a corresponding standard free energy

\begin{equation}
    \Delta \textrm{G}^\circ_{\textrm{non-covalent}} = \textrm{-RT} \ln \textrm{K}_\textrm{I}
\label{eq:covbind_eq2b}
\end{equation}

Thus, all of the concepts developed for non-covalent binding (enthalpy/entropy balance, solvation, conformational and translational/rotational entropy) apply directly to this pre-equilibrium\cite{CH3_Strelow2017}.

The second step is a chemical reaction in which a nucleophilic residue on the protein (e.g. Cys, Ser, Lys) attacks an electrophilic warhead on the ligand. Its rate is governed by the inactivation rate constant $\textrm{k}_{\textrm{inact}}$, which is determined by an activation free energy $\Delta G^\ddagger$ using the transition state theory:\cite{CH3_Schramm2011}

\begin{equation}
    \textrm{k}_{\textrm{inact}} \approx \frac{k_B T}{h} \exp\!\left(-\frac{\Delta G^\ddagger}{RT}\right)
\label{eq:covbind_eq3}
\end{equation}

with $\Delta G^\ddagger = \Delta H^\ddagger - T \Delta S^\ddagger$. The overall efficiency of a covalent inhibitor is then quantified by the second-order parameter $k_{\text{inact}}/K_I$, which combines the thermodynamics of noncovalent complex formation with the thermodynamics of the chemical step\cite{CH3_Strelow2017}. In addition, the reaction free energy $\Delta G^\circ_{\text{rxn}}$ for forming the covalent adduct is often strongly negative, making dissociation negligible on experimental timescales and conferring prolonged target occupancy characteristic of covalent drugs.

\section{Scoring functions}
Docking workflows use scoring functions for pose prediction (finding bioactive conformation), affinity prediction (estimating free energy of binding), and virtual screening (finding new hit molecules from large databases). Since scoring functions are developed using experimental binding data along with structural data from the protein data bank, their performance, such as pose prediction, scoring, and ranking, must be validated using external benchmark sets that were not a part of the dataset for development\cite{CH3_Su2019_CASF2016,CH3_Wang2014_CASF2013_Set,CH3_Wang2014_CASF2013_Results}.

Here, we focus on the classes of scoring functions used in docking, highlighting their mathematical forms and underlying physical/statistical assumptions. We also discuss parameterisation, testing, and validation of scoring functions. For validation, we enlist and discuss various datasets and benchmarks (e.g., PDBbind, CASF, DUD-E, CSAR) that have emerged over the years. These benchmark sets and datasets have been curated with the help of the molecular modelling community, which conducts blind challenges to rigorously test every aspect of a molecular docking program's capabilities. We also highlight key challenges, including data bias, overfitting, protein flexibility, solvation, and generalisation, which motivates current research for improving the docking methods in general, and scoring functions in particular\cite{CH3_Wang2017} 

\section{Classes of Scoring Functions}
 Protein--ligand scoring functions are classified under four classes based on their parameterisation philosophy, and are given as follows:

\begin{enumerate}
    \item Force field-based (physics-based) scoring functions.
    \item Empirical scoring functions.
    \item Knowledge-based (statistical potential) scoring functions.
    \item Machine learning-based scoring functions.
\end{enumerate}

In recent times, several good scoring functions in use today are consensus scoring functions that combine elements from different scoring functions. They are further fine-tuned using suitable weighting functions or fitting coefficients to reproduce biological binding affinity. Nevertheless, for historical reasons, we explain each class of scoring functions and their philosophies. This should prepare the reader to appreciate the modern scoring functions. 

\subsection{Force field-based scoring functions}
Force field-based scoring functions are derived directly from classical molecular mechanics (MM) energy expressions. A typical MM energy for a protein--ligand complex is written as:

\begin{equation}
    \textrm{E}_{\textrm{MM}} = \textrm{E}_{\textrm{bond}} + \textrm{E}_{\textrm{angle}} + \textrm{E}_{\textrm{tors}} + \textrm{E}_{\textrm{vdW}} + \textrm{E}_{\textrm{elec}}
\label{eq:FF_general}
\end{equation}

where the bonded terms model internal strain, and the nonbonded terms account for intermolecular interactions, further classified as van der Waals (vdW) and electrostatic interactions. In docking, the bonded terms within the protein are often held fixed (rigid receptor assumption). In contrast, bonded terms within the ligand and cross terms (e.g., ligand dihedrals coupled to receptor) may be partially considered. However, for most practical reasons, components describing bonds and angles are held fixed for both ligand and receptors. 

A generic form of the nonbonded MM energy between protein and ligand can be written as
\begin{equation}
    \textrm{E}_{\textrm{PL}}^{\textrm{MM}} =
    \sum_{i  \in  \textrm{protein}} \sum_{j  \in  \textrm{ligand}}
    \left[
    \frac{A_{ij}}{r_{ij}^{12}}
    -
    \frac{B_{ij}}{r_{ij}^{6}}
    +
    \frac{q_i q_j}{4\pi \varepsilon_0 \varepsilon_r r_{ij}}
    \right],
\label{eq:MM_nonbonded}
\end{equation}

where $\textrm{r}_{\textrm{ij}}$ is the distance between ligand atoms i and receptor atoms j, $\textrm{A}_{\textrm{ij}}$ and $\textrm{B}_{\textrm{ij}}$ are Lennard--Jones repulsive and attractive parameters respectively, $\textrm{q}_\textrm{i}$ and $\textrm{q}_\textrm{j}$ are partial charges, and $\varepsilon_r$ represents an effective dielectric constant. The Lennard-Jones term approximates dispersion and steric repulsion, while the Coulombic term captures electrostatics. 

In a purely force field-based scoring function, the docking score is often taken to be
\begin{equation}
    \textrm{S}_{\textrm{FF}} = \textrm{E}_{\textrm{PL}}^{\textrm{MM}} + \textrm{E}_{\textrm{env}}
\label{eq:MM_scoring}
\end{equation}

where $\textrm{E}_{\textrm{env}}$ is an environmental term that approximately accounts for solvation and sometimes entropic contributions. In docking implementations, $\textrm{E}_{\textrm{env}}$ is frequently modelled by implicit solvent schemes, such as generalised Born (GB) or Poisson--Boltzmann (PB) plus a surface area term, leading to MM/GBSA or MM/PBSA-type rescoring\cite{CH3_wang2019end}. 

Even when not explicitly labelled as force field-based, many docking programs employ grid-based precomputation of the receptor's nonbonded potential for each atom type, effectively discretising Equation \ref{eq:MM_nonbonded} on a 3D lattice. The ligand is then placed in this field, and its energy is interpolated for efficiency.

The strengths of force field-based scoring functions include:
\begin{itemize}
    \item Clear physical interpretation of terms.
    \item Transferability across chemical groups when parameters are well defined.
    \item More amenable to rigorous free-energy methods (FEP, TI) for post-processing.
\end{itemize}

And the limitations include:
\begin{itemize}
    \item Difficulty in accurately modelling desolvation, and entropic contributions in a simple additive energy function.
    \item Sensitivity to protonation states, tautomers, and partial charge assignments.
\end{itemize}

\subsection{Empirical scoring functions}
Empirical scoring functions approximate the binding free energy as a weighted sum of physics-based descriptors fitted against experimental binding data. 

A general functional form is given as:
\begin{equation}
    \Delta \textrm{G}_{\text{bind}}^{\text{emp}} = w_0 + \sum_{k=1}^{K} w_k f_k(\mathbf{R})
\label{eq:empirical_general}
\end{equation}

where R denotes the protein--ligand complex, $\textrm{f}_\textrm{k}(\textrm{R})$ are descriptor functions (e.g., counts of hydrogen bonds, hydrophobic contacts, metal coordination interactions, buried surface area, number of rotatable bonds), and $w_k$ are wieghting functions or coefficients fitted by regression to experimental binding free energies (converted from $\textrm{K}_\textrm{d}$, $\textrm{K}_\textrm{i}$, or IC$_{\textrm{50}}$ values)\cite{CH3_Wang2017}.

For example, classical empirical scoring functions such as ChemScore and X-Score use terms like
\begin{equation}
    \Delta \textrm{G}_{\textrm{bind}} = \Delta \textrm{G}_{\textrm{vdW}} + \Delta \textrm{G}_{\textrm{Hbond}} + \Delta \textrm{G}_{\textrm{lipophilic}} + \Delta \textrm{G}_{\textrm{rot}} + \Delta \textrm{G}_\textrm{0}
\label{eq:classicalEMP}
\end{equation}

where each $\Delta \textrm{G}$ term is itself a linear function of structural descriptors (e.g., hydrogen bond counts weighted by geometry-dependent factors), these parameters are obtained by least-squares fitting to the training set of complexes with high-quality structures and experimentally measured binding affinities\cite{CH3_Wang2017}.

Empirical scoring functions are generally fast to evaluate, making them an attractive choice for high-throughput virtual screening. 

Their limitations include:
\begin{itemize}
    \item Dependence on the composition and quality of the training set, which may introduce several biases.
    \item Linear additivity assumptions may fail in the presence of cooperative effects (e.g., multiple hydrogen bonds in a constrained geometry).
    \item Limited explicit treatment of water-mediated interactions, protonation equilibria, and protein flexibility.
\end{itemize}

\subsection{Knowledge-based scoring functions}

Knowledge-based scoring functions, also called statistical potentials or potentials of mean force (PMFs), derive effective interaction potentials from the observed contact frequencies in a structural database, typically the Protein Data Bank (PDB) or curated subsets such as PDBbind\cite{CH3_Wang2017}.  The central idea is the inverse Boltzmann relation, which states that \emph{if a particular atom-type pair at a certain distance occurs more (or less) frequently than expected by chance, then that configuration is associated with an effective favourable (or unfavourable) free energy.}

Let $\alpha$ and $\beta$ denote atom types (e.g., protein carbonyl oxygen and ligand aromatic carbon). The radial distribution function $g_{\alpha\beta}(r)$ is estimated from the training set by counting occurrences of atom-type pairs at distance r, and normalising by a reference distribution $g_{\alpha\beta}^{\text{ref}}(r)$ (often the distribution for a non-interacting reference state). 

The potential of mean force is then calculated as:

\begin{equation}
    \textrm{U}_{\alpha\beta}(r) = - k_B T \, \ln \left( \frac{g_{\alpha\beta}(r)}{g_{\alpha\beta}^{\text{ref}}(r)} \right),
\label{eq:PMF1}
\end{equation}

where $\textrm{k}_\textrm{B}$ is Boltzmann's constant and T is the temperature.

Given a protein--ligand complex, the total score is computed as a sum over all intermolecular atom pairs:
\begin{equation}
    \textrm{S}_{\textrm{PMF}} = \sum_{i \in \textrm{protein}} \sum_{j \in \textrm{ligand}} \textrm{U}_{\alpha(i)\beta(j)}\bigl(r_{ij}\bigr)
\label{eq:PMF2}
\end{equation}
where $\alpha(i)$ and $\beta(j)$ denote atom-type from the protein (i) and ligand (j).
Such statistical potentials can approximate complex many-body effects, such as solvation, packing, and indirect interactions, implicitly encoded in the structural database. These lead to several advantages such as:

\begin{itemize}
    \item Incorporation of large-scale structural information from many protein--ligand complexes.
    \item Straightforward extension to new atom types or interaction classes by recomputing statistical weights.
\end{itemize}

However, the main observed disadvantage is their strong dependence on the choice of reference state and on the representativeness of the training dataset. Furthermore, they may not extrapolate well beyond the domain of the observed chemical class. This may be remedied by either continuous development to incorporate new atom-pair information or on-the-fly learning, which is mainly observed in machine-learning or deep-learning based scoring functions.

\subsection{Consensus Scoring functions}

Consensus scoring functions emerged because no single classical scoring function consistently performs reliably across all targets, chemotypes, and docking scenarios. Force field, empirical, and knowledge-based scoring functions each approximate binding thermodynamics differently, with unique strengths and systematic limitations. Consensus strategies integrate these approaches to leverage their partial complementarity, resulting in more robust pose ranking and virtual screening performance than any individual score\cite{CH3_Wang2017,CH3_nhat2023}. To combine different scoring functions to develop a consensus scoring functions, following three methods are used to combine other scoring functions and assign an appropriate weight to them:

\begin{itemize}
    \item Score-based consensus methods.
    \item Rank-based consensus methods.
    \item Voting or classifier-style consensus methods.
\end{itemize}   
Score-based consensus methods\cite{CH3_yang2005}, the simplest category, combine numeric scores from multiple functions directly. After normalisation, such as z-score or min–max scaling to account for differing score ranges and directions, an unweighted or weighted average of scores from force field, empirical, and knowledge-based functions can be calculated. Weighted schemes adjust each function's contribution to maximize correlation with experimental data on a training or validation set; for example, a more reliable empirical function may receive greater weight than a force field term. Score-based consensus methods are attractive due to their straightforward implementation and their ability to maintain a continuous score scale suitable for rank ordering and thresholding.

Rank-based consensus methods ignore score magnitudes and instead combine the rank positions assigned by each scoring function. For example, Charifson and co-workers\cite{CH3_charifson1999} docked compound libraries using several empirical and force field-derived scores, then selected compounds within the top fraction (e.g., the top 10–20\%) based on multiple functions, effectively intersecting high-ranking subsets. This approach improved hit rates and reduced false positives compared to any individual scoring function, demonstrating that agreement among classical scores is a valuable heuristic for enrichment. Additional rank-based schemes, such as Borda counts or average ranks\cite{CH3_grillberger2026}, are less stringent than strict intersection but still prioritise ligands consistently ranked highly across empirical, knowledge-based, and force-field scores.

The third category includes voting or classifier-style consensus methods\cite{CH3_votedock2011}, in which each scoring function votes on the likelihood that a pose or ligand is active, and the final decision is determined by majority or unanimity rules. For example, in structure-based virtual screening (see \textbf{Equation \ref{eq:conscore}}), a ligand may be classified as promising if two out of three classical scores exceed target-specific thresholds. This approach provides a coarse but pragmatic filter for prioritising hits before more computationally intensive rescoring or experimental validation.

\begin{equation}
\begin{aligned}
\text{MetaScore} = -0.358 \times \text{HTScore} + 0.015 \times \text{FlexXScore} \\\\ 
-0.358 \times \text{GlideScore} + 0.014 \times \text{GoldScore} \\\\
+0.004 \times \text{LigScore} + 0.15 \times \text{SurflexScore}
\end{aligned}
\label{eq:conscore}
\end{equation}

Numerous retrospective studies have shown that consensus scoring, constructed from combinations of empirical, knowledge-based, and force-field-based functions, generally yields more stable enrichment factors and poses prediction success rates than any individual score\cite{CH3_konstantinou2007} However, consensus scoring does not address fundamental limitations shared by its components, such as restricted protein flexibility and simplistic solvation models, and performance improvements remain target dependent.


\subsection{Machine-Learning-Based Scoring Functions}

Machine-learning (ML) scoring functions replace manually fitted linear or pairwise potentials with more flexible models trained to reproduce binding data. Early ML scoring functions used relatively simple descriptors and regression algorithms (e.g., random forests), while recent work employs deep neural networks operating directly on 3D grids or molecular graphs\cite{CH3_Ballester2010_RFScore,CH3_Wojcikowski2017_RFScoreVS,CH3_Jimenez2018_KDEEP}.

In a generic ML scoring function, the goal is to learn a mapping.
\begin{equation}
    \hat{y} = f_{\boldsymbol{\theta}}(\mathbf{x}),
\end{equation}

where $\mathbf{x} \in \mathbb{R}^d$ encodes features of the protein--ligand complex (e.g., atom-type pair counts, interaction fingerprints, grid-based densities), $y$ is the target property (e.g., $\Delta G_{\text{bind}}$ or $pK_d$), and $\boldsymbol{\theta}$ are model parameters. The parameters are optimised by minimising a loss function over a training set $\{(\mathbf{x}_n, y_n)\}_{n=1}^N$, typically the mean-squared error.

\begin{equation}
    \mathcal{L}(\boldsymbol{\theta}) = \frac{1}{N} \sum_{n=1}^N \left( y_n - f_{\boldsymbol{\theta}}(\mathbf{x}_n) \right)^2.
\end{equation}

\subsubsection{Feature-based ML scoring functions}
RF-Score is a prototypical feature-based ML scoring function\cite{CH3_Ballester2010_RFScore}. It uses a fixed-length descriptor vector whose components are counts of protein--ligand atom-type pairs within a cutoff radius (e.g., 12 \r{A}). For example, one feature may count all occurrences where a protein carbon atom lies within 6--8 \r{A} of a ligand nitrogen atom. A random forest of decision trees is trained to predict binding affinity from these descriptors. The RF-Score prediction is an average over M trees,

\begin{equation}
    \hat{y} = \frac{1}{M} \sum_{m=1}^{M} T_m(\mathbf{x}),
\end{equation}
where $T_m$ denotes the prediction from the $m$-th tree.

RF-Score-VS\cite{CH3_Wojcikowski2017_RFScoreVS} extends this idea to virtual screening tasks, training on large numbers of active/decoy complexes (e.g., DUD-E) and optimising enrichment metrics. Related work\cite{CH3_li2021machine,CH3_li2017structural} has shown that ML scoring functions can substantially improve virtual screening hit rates relative to classical scoring functions, but are sensitive to training/test similarity and dataset biases.

\subsubsection{Deep learning scoring functions}
Deep learning methods operate directly on spatial or graph representations of protein--ligand complexes. In 3D convolutional neural network (3D-CNN) approaches, atoms are mapped to voxelized grids with channels corresponding to atom types or physicochemical properties, and convolutions learn spatial patterns predictive of binding affinity\cite{CH3_Jimenez2018_KDEEP}. Other models use graph neural networks (GNNs) where nodes represent atoms and edges represent bonds or noncovalent contacts; message-passing layers propagate information and construct learned representations of the complex\cite{CH3_Wang2017}.

While these models often achieve state-of-the-art performance on standard benchmarks such as CASF-2013 and CASF-2016, recent work has emphasised that their apparent gains can be inflated by data leakage and training--test similarity in PDBbind-derived splits.\cite{CH3_li2021machine,CH3_li2017structural} Addressing these issues requires careful dataset design and evaluation, discussed further below.

\section{Datasets and validation for Scoring Functions}
Reliable development and evaluation of scoring functions require curated datasets that link high-quality structural data on protein--ligand complexes with experimentally measured binding affinity data. The following sections briefly discuss commonly used, publicly available datasets. 

\subsection{PDBbind Database}
PDBbind was created to provide a comprehensive collection of binding affinities for biomolecular complexes in the Protein Data Bank (PDB).\cite{CH3_Wang2004_PDBbind,CH3_Wang2005_PDBbind,CH3_Liu2015_PDBbind}. It links PDB entries to experimentally measured $\textrm{K}_\textrm{d}$, $\textrm{K}_\textrm{i}$, or IC$_{\textrm{50}}$ values from the primary literature and provides processed structural files suitable for docking and scoring studies.

PDBbind is organised into three key subsets:

\begin{itemize}
    \item \textbf{General set}: All complexes with available binding data and acceptable structural quality.
    \item \textbf{Refined set}: A high-quality subset filtered by criteria such as resolution ($\leq 2.5$ \r{A}), absence of severe structural issues, and reliable binding constants.
    \item \textbf{Core set}: A non-redundant subset clustered by protein sequence identity (e.g., 90\%) with a small number of representatives per cluster. This core set is used as the test set in the CASF benchmarks.
\end{itemize}

Generally, classical scoring functions have been trained or evaluated using the PDBbind dataset\cite{CH3_Trott2010_Vina,CH3_Wang2017}. However, PDBbind is not free from biases and artefacts: redundancy, uneven coverage of target classes, and issues with protonation states or ligand modelling can all affect performance estimates. Recent developments have proposed reorganised splits (e.g., structure-aware or sequence-dissimilar splits) and leak-proof variants that reduce train--test overlap\cite{CH3_li2017structural,CH3_li2021machine}.

\subsection{CASF Benchmarks} \label{sec:CASF}
The Comparative Assessment of Scoring Functions (CASF) benchmark series was developed to provide standardised, objective evaluations of scoring functions\cite{CH3_Wang2014_CASF2013_Set,CH3_Wang2014_CASF2013_Results,CH3_Su2019_CASF2016}. CASF decouples scoring from docking by predefined sets of protein--ligand poses, usually near-native plus decoys, so that different scoring functions can be compared under identical sampling conditions.

In CASF-2013, a core set of 195 high-quality complexes from PDBbind v2013 was assembled\cite{CH3_Wang2014_CASF2013_Set}. CASF-2016 updated this to 285 complexes and improved evaluation protocols\cite{CH3_Su2019_CASF2016}.

For each scoring function, four types of performance are assessed:

\begin{enumerate}
    \item \textbf{Scoring power}: Correlation (Pearson $R_p$) between predicted scores and experimental binding affinities across the test set.
    \item \textbf{Ranking power}: Ability to correctly rank ligands binding to the same target by affinity.
    \item \textbf{Docking power}: Ability to identify near-native poses (e.g., RMSD $\leq 2.0$ \r{A}) among decoy poses.
    \item \textbf{Screening power}: Enrichment of true actives spiked in the set of decoys in a virtual screening scenario.
\end{enumerate}

These benchmarks have become the \textit{de facto} standard for reporting the performance of new scoring functions.

\subsection{DUD and DUD-E for Virtual Screening}
For virtual screening, datasets of actives and carefully designed decoys are required. The Directory of Useful Decoys (DUD) and its successor DUD-E provide such benchmarks\cite{CH3_Mysinger2012_DUDE}. DUD-E contains 102 protein targets with 22,886 active ligands and, for each active, 50 property-matched decoys drawn from the ZINC database. Decoys are chosen to match key physicochemical properties (e.g., molecular weight, cLogP, number of hydrogen-bond donors/acceptors) while being topologically dissimilar, reducing trivial discrimination by simple property filters. DUD-E is widely used to assess the screening power of docking and scoring methods using metrics such as the enrichment factor (EF), the receiver-operating characteristic (ROC) area under the curve (AUC), and the Boltzmann-enhanced discrimination of ROC (BEDROC).

\subsection{CSAR Benchmarks}
The Community Structure--Activity Resource (CSAR) benchmark exercises have provided a series of blind and retrospective challenges for docking, scoring, and ranking\cite{CH3_CSAR2013}.

CSAR benchmark sets include:
\begin{itemize}
    \item Pose prediction tasks with pre-generated decoy poses for multiple ligands.
    \item Affinity ranking tasks where participants must rank ligands by experimental binding data.
    \item Decoupled scoring tasks where only scoring (not docking) is evaluated by providing fixed pose ensembles.
\end{itemize}

These exercises emphasise realistic, prospectively challenging problems and highlight limitations of current methods, particularly in handling protein flexibility, water networks, and subtle structure--activity relationships.

\subsection{Other Datasets}
Additional datasets used to develop and validate scoring functions include:

\begin{itemize}
    \item \textbf{Astex Diverse Set}\cite{CH3_hartshorn2007}: A curated set of protein--ligand complexes with high-quality crystal structures, often used to test docking pose prediction.
    \item \textbf{sc-PDB}\cite{CH3_scPDB2014}: A database of ligandable binding sites extracted from the PDB, designed for docking and pharmacophore modelling.
    \item \textbf{DEKOIS 2.0}\cite{CH3_bauer2013}: A collection of challenging benchmarks for virtual screening, complementing DUD-E by focusing on difficult targets and well-matched decoys.
    \item \textbf{Cross-docked sets}: Large datasets where ligands are redocked into multiple related receptor structures to assess robustness to receptor conformational variation.
\end{itemize}

Together with PDBbind and CASF, these resources form the backbone of modern scoring function development and evaluation.

\section{Development of Scoring Functions}

This section outlines a typical workflow for developing protein--ligand scoring functions, from problem definition and data curation to model form selection and parameter fitting. The process differs in detail between classical (empirical/knowledge-based/force-field) and ML-based approaches, but shares many common elements\cite{CH3_Wang2017}.

\subsection{Problem Definition and Target Application}
Scoring functions are rarely universal; they are implicitly optimised for specific tasks. Three main use cases are:

\begin{itemize}
    \item \textbf{Pose prediction}: Selecting the near native pose or bioactive conformations of a ligand bound to a known receptor structure.
    \item \textbf{Binding affinity prediction}: Predicting absolute or relative binding free energies across diverse ligands and targets.
    \item \textbf{Virtual screening}: Discriminating actives from inactives in large ligand libraries.
\end{itemize}

For example, AutoDock-Vina's scoring function was trained primarily to improve docking (pose prediction) performance relative to AutoDock4, rather than to maximise correlation with experimental binding affinities\cite{CH3_Trott2010_Vina}.

\subsection{Data Curation and Preprocessing}
High-quality data are essential for meaningful scoring function development\cite{CH3_Wang2005_PDBbind,CH3_Liu2015_PDBbind}. Typical pre-processing steps include:

\begin{enumerate}
    \item \textbf{Structure selection}: Choose complexes with good resolution (e.g., $\leq 2.5$ \r{A}), low R-factors, no severe clashes or missing residues in the binding site, and well-resolved ligands.
    \item \textbf{Binding data extraction and harmonisation}: Extract $K_d$, $K_i$, or IC$_{50}$ values from primary literature, ensuring that the reported measurements correspond to the crystallised complex and similar experimental conditions.
    \item \textbf{Conversion to uniform targets}: Convert binding constants to binding free energies using the equation \ref{eq:dg_ka}
    \item \textbf{Structure preparation}: Add missing hydrogens, resolve alternate conformations, optimise side-chain orientations in the binding site if necessary, and standardise atom naming and types.
    \item \textbf{Protonation and tautomer states}: Assign physiologically appropriate protonation states and tautomers at a reference pH (typically $\sim 7$), as these strongly influence electrostatics and hydrogen bonding.
    \item \textbf{Redundancy control}: Cluster complexes by sequence similarity and/or ligand similarity to control redundancy and avoid near-duplicate complexes in training and test sets.
\end{enumerate}

\subsection{Descriptor Design and Model Form}

The choice of descriptors $f_k(\mathbf{R})$ and model form $f_{\boldsymbol{\theta}}$ distinguishes different classes of scoring functions.


\subsubsection{Classical (Empirical, Force-Field, Knowledge-Based) Models}

Classical scoring functions define descriptors with clear physics-based or structural interpretations:

\begin{itemize}
    \item Number and geometry of hydrogen bonds.
    \item Hydrophobic contact surface area or counts of hydrophobic atom contacts.
    \item Metal coordination interactions.
    \item Buried solvent-accessible surface area or polar surface area changes.
    \item Ligand flexibility (e.g., number of rotatable bonds).
    \item Pairwise atom-type contact frequencies (for PMFs).
\end{itemize}

These descriptors feed into relatively simple functional forms, such as linear combinations (\textbf{Equation \ref{eq:empirical_general}}) or pairwise additive potentials (\textbf{Equation \ref{eq:PMF2}}).

\subsection{Parameter Estimation and Training}

In empirical models, the parameters $w_k$ are typically estimated by linear or non-linear regression. A common approach is to minimise the sum of squared differences between predicted and experimental $\Delta G_{\text{bind}}$ over a training set:

\begin{equation}
    \min_{\{w_k\}} \sum_{n=1}^{N} \left( \Delta G_{\text{bind},n}^{\text{exp}} - \Delta G_{\text{bind},n}^{\text{emp}} \right)^2
\label{eq:parameter_fitting}
\end{equation}

 Constraints or normalisation may be applied to keep parameters within the physics-based principles.
AutoDock Vina adopted a more complex optimisation strategy, tuning not only linear weights but also functional forms (e.g., Gaussian potential widths) using PDBbind as a training set and optimising performance across both scoring and docking tasks\cite{CH3_Trott2010_Vina}.

Knowledge-based potentials rely less on explicit parameter fitting and more on statistical estimation of $g_{\alpha\beta}(r)$, and the choice of reference state. Some smoothing or normalisation is usually applied to avoid overfitting.

\section{Testing and Validation of Scoring Functions}
Rigorous validation is essential to assess how well a scoring function will perform on new, unseen targets and ligands.

\subsection{Internal Validation: Cross-Validation and Bootstrapping}
During development, internal validation methods such as k-fold cross-validation or bootstrapping are often used:

\begin{itemize}
    \item In k-fold cross-validation, the dataset is split into k subsets; each subset is held out in turn as a test set while the model is trained on the remaining k-1 subsets.
    \item Bootstrapping involves sampling (with replacement) multiple training sets and evaluating performance on the out-of-bag samples.
\end{itemize}

These methods provide estimates of predictive performance; however, they can be misleading if random splits result in highly similar complexes appearing in both the training and test sets. Structure-aware splitting schemes that enforce dissimilarity in protein sequence and/or ligand scaffolds are now recommended for ML scoring functions\cite{CH3_li2017structural}.

\subsection{External Benchmarks: CASF}

As already discussed above (see section \ref{sec:CASF}), the CASF benchmarks provide a standardised external evaluation of scoring functions\cite{CH3_Wang2014_CASF2013_Results,CH3_Su2019_CASF2016}.
Key metrics include:
\begin{itemize}
    \item \textbf{Scoring power}: Pearson correlation coefficient $\textrm{R}_\textrm{p}$ between predicted and experimental binding affinities ($\Delta \textrm{G}_{\textrm{bind}}$ or $\textrm{pK}_\textrm{d}$) over the core set. Root-mean-square error (RMSE) is also reported.
    \item \textbf{Ranking power}: Percentage of correctly ranked pairs or triplets of ligands binding to the same target; rank correlation coefficients (Spearman $\rho$, Kendall $\tau$) are used.
    \item \textbf{Docking power}: Percentage of complexes for which the top-ranked pose is within a specified RMSD threshold (e.g., 2.0 \r{A}) of the crystal pose.
    \item \textbf{Screening power}: Enrichment factors or ROC AUC when discriminating actives from decoys.
\end{itemize}

Comparative CASF studies have shown that classical scoring functions tend to perform better in docking and virtual screening than in absolute scoring and ranking, reflecting the difficulty of accurately capturing subtle free-energy differences\cite{CH3_Su2019_CASF2016,CH3_Wang2014_CASF2013_Results}.

\subsection{Virtual Screening Benchmarks: DUD-E and Others}
To evaluate virtual screening performance, scoring functions are applied to datasets such as DUD-E\cite{CH3_Mysinger2012_DUDE}, DEKOIS-2.0\cite{CH3_bauer2013}, and DUDE-Z.\cite{CH3_Mysinger2012_DUDE}
Common metrics include:
\begin{itemize}
    \item \textbf{Enrichment factor}: The ratio of the fraction of actives retrieved in the top-ranked set to the fraction expected by random selection.
    \item \textbf{ROC AUC}: Area under the ROC curve, measuring discrimination across all thresholds.
    \item \textbf{BEDROC}: Emphasizes early recognition of actives.
\end{itemize}

\subsection{Pose prediction and CSAR benchmarks}
Pose prediction benchmarks, including those in CSAR exercises, provide decoy pose sets for which only one or a few near-native (close to bioactive or experimentally observed) poses exist\cite{CH3_CSAR2013}. These tests directly assess the ability of a score to recognise the correct binding pose given an ensemble of candidate poses. They also highlight the interplay between sampling and scoring: even an excellent scoring function cannot succeed when the correct pose is absent from the candidate set.

Scoring functions are evaluated by:
\begin{itemize}
    \item Percentage of systems where a near-native pose is ranked in the top N.
    \item Distributions of RMSD for top-ranked poses.
\end{itemize}

\section{Current Challenges}
Despite years of development, current scoring functions still fall short of fully reliable binding affinity prediction, especially for prospective applications\cite{CH3_Wang2017,CH3_Su2019_CASF2016}. Important challenges include:

\begin{itemize}
    \item \textbf{Desolvation and water-mediated interactions}: Accurately modelling the energetic cost of desolvating both ligand and binding site, as well as the contributions of displaced, retained, or bridging water molecules, remains difficult in simple additive scoring functions.
    \item \textbf{Protein flexibility}: Most scoring functions assume a rigid or semi-rigid receptor; large conformational changes upon binding are poorly captured.
    \item \textbf{Entropy}: Configurational and conformational entropies of ligand, receptor, and solvent are only crudely approximated (e.g., via penalties for rotatable bonds).
    \item \textbf{Protonation and tautomerism}: Binding often involves shifts in protonation states or tautomeric forms; static protonation assignments can be incorrect and largely misleading.
    \item \textbf{Transferability}: Parameters optimised on limited datasets may not be transferable to new targets, especially those with different atom-types or atom-pairs.
\end{itemize}

Consensus approaches that combine physics-based methods (e.g., MM/GBSA, FEP) with ML corrections, as well as more physics-informed ML architectures (e.g., equivariant neural networks) that respect symmetries and conservation laws, are promising avenues. There is also growing emphasis on uncertainty quantification and domain-of-applicability estimation, so users can gauge confidence in scoring predictions rather than relying solely on point estimates. On the data side, improved curation of structural and binding data, careful splitting strategies, and new high-quality datasets (including prospective validation campaigns) will be critical to drive progress\cite{CH3_Liu2015_PDBbind}.

\section{Summary}
Scoring functions are central to molecular docking and structure-based drug design, providing a bridge between 3D molecular structures and quantitative assessments of binding. Classical force field, empirical, and knowledge-based scoring functions have laid the foundation for modern docking workflows. At the same time, machine learning-based models have demonstrated tremendous potential for substantial improvements supported by high-quality data and rigorous validation. However, no single approach is universally superior: each scoring function reflects trade-offs among physical interpretability, computational efficiency, and flexibility to learn complex structure — affinity relationships. 

Continued progress will likely rely on:
\begin{itemize}
    \item Better integration of physics-informed machine-learning methods.
    \item More rigorous benchmarking frameworks that minimise leakage and bias.
    \item Prospective validation of scoring functions in real drug discovery campaigns.
\end{itemize}

\bibliographystyle{unsrtnat}
\bibliography{chapter3}

\end{document}